\section*{Additional abbreviations introduced in P1E--P1F}
\begin{longtable}{>{\bfseries}p{0.20\textwidth}p{0.72\textwidth}}
\toprule
Abbreviation & Meaning\\
\midrule
\endfirsthead
\toprule
Abbreviation & Meaning\\
\midrule
\endhead
MAP & Maximum A Posteriori; in the particle implementation, the state or lineage associated with the particle having the largest pre-transition posterior weight.\\
ZUPT & Zero-Velocity Update; an inertial-navigation correction used in the pedestrian processing chain discussed by Han et al.\\
\bottomrule
\end{longtable}

\section*{P1E--P1F paper-reproduction quantities}
The symbols in this table belong to the separate three-state reproduction of Han et al. and should not be confused with the five-state IMU-driven thesis estimator defined above.
\begin{longtable}{p{0.30\textwidth}p{0.62\textwidth}}
\toprule
Symbol & Meaning\\
\midrule
\endfirsthead
\toprule
Symbol & Meaning\\
\midrule
\endhead
$\bm x_t^{\mathrm{paper}}$ & Paper-level target state $[x_t,y_t,\phi_t]^\top$ used in P1E--P1F. The superscript ``paper'' is used in this notation table only to distinguish the model from the five-state thesis baseline; equations may omit it when the context is unambiguous.\\
$x_t$ & Global target $x$-coordinate in the paper-level navigation frame at sample $t$.\\
$y_t$ & Global target $y$-coordinate in the paper-level navigation frame at sample $t$.\\
$\phi_t$ & Global target azimuth/yaw in the paper-level model.\\
$\bm u_t^{\mathrm{paper}}$ & Paper-level DR input $[\Delta L_t,\Delta\phi_t]^\top$.\\
$\Delta L_t$ & Travelled-distance increment supplied by the higher-level DR system for one propagation interval.\\
$\Delta\phi_t$ & Azimuth increment supplied by the higher-level DR system for one propagation interval.\\
$N$ & Number of initial position hypotheses in the $N\times M$ notation used by Han et al.\\
$M$ & Number of initial azimuth hypotheses associated with the paper's $N\times M$ particle construction.\\
$N_p=NM$ & Total number of particles in the paper-level reproduction when the $N\times M$ notation is used.\\
$\theta_i$ & Position-bearing angle used to place initial particle $i$ around the known auxiliary node.\\
$r_i$ & Initial radial distance from the auxiliary node to particle $i$.\\
$\delta r_i$ & Radial perturbation applied to the first measured UWB range when constructing particle $i$.\\
$\Delta d$ & Half-width of the first-range annulus. P1F uses the repository choice $\Delta d=3\sigma_{\mathrm{UWB}}$.\\
$q(\bm x_t^i\mid\bm x_{t-1}^i,z_t)$ & Generic particle-filter proposal density appearing in the source formulation.\\
$\ell_t^i$ & Log-likelihood assigned to particle $i$ at sample $t$ from the UWB range residuals.\\
$\rho_j$ & Range residual associated with auxiliary node $j$ in the adaptive auxiliary-rejection mechanism of Han et al.\\
$\Omega_1,\Omega_2$ & Source-defined but numerically unspecified thresholds for residual mean and residual variance in adaptive auxiliary rejection.\\
$\mathcal C_i$ & Set of candidate destination particles considered for ACO movement of source particle $i$.\\
$\Delta w_{ij}$ & Weight-information difference between candidate particle $j$ and source particle $i$; the repository interpretation is $w_j-w_i$ for $w_j>w_i$.\\
$\eta_{ij}$ & Position-information factor between particles $i$ and $j$; the repository interpretation is inverse Euclidean distance.\\
$\epsilon_d$ & Small positive distance regularizer used to avoid division by zero in $\eta_{ij}$.\\
$\epsilon_w$ & Small positive weight regularizer used in the ACO transition score.\\
$\alpha$ & Exponent controlling the contribution of weight information to the ACO transition score.\\
$\beta$ & Exponent controlling the contribution of spatial-distance information to the ACO transition score.\\
$s_{ij}$ & Unnormalized ACO transition score from source particle $i$ to candidate particle $j$.\\
$P_{ij}$ & Normalized ACO transition probability/score fraction from source particle $i$ to candidate particle $j$.\\
$j^*$ & Candidate index with the maximum transition probability for a given source particle.\\
$\lambda$ & ACO movement threshold: movement occurs only when the maximum candidate transition probability exceeds this value.\\
$K_i$ & Number of candidate destination particles in $\mathcal C_i$.\\
$c_\lambda$ & Candidate-count-normalized transition-threshold factor defined by $c_\lambda=\lambda K_i$. In P1F-B the implementation equivalently uses $\lambda_i=c_\lambda/K_i$.\\
$s_k^{\mathrm{paper}}$ & Deterministic speed profile used only to generate the P1F-A paper-state truth trajectory.\\
$\omega_k^{\mathrm{paper}}$ & Deterministic azimuth-rate profile used only to generate the P1F-A paper-state truth trajectory.\\
$e_{p,k}^{(i)}$ & Position error of paper-state particle $i$ relative to the known simulation truth at sample $k$.\\
$e_{\phi,k}^{(i)}$ & Wrapped azimuth error of paper-state particle $i$ relative to the known simulation truth.\\
$W_{\mathrm{mode},k}$ & Posterior probability mass contained in the diagnostic correct-mode region at sample $k$, evaluated before ACO movement in P1F-B.\\
$C_{\mathrm{mode},0}$ & Number of initial particles lying in the diagnostic correct-mode region.\\
$\phi_{0,\mathrm{MAP},k}$ & Initial-azimuth lineage carried by the MAP particle at sample $k$.\\
$N_{\mathrm{move},k}$ & Number of source particles that copy a destination during the ACO transition at sample $k$.\\
$f_{\mathrm{move},k}$ & Moved-particle fraction $N_{\mathrm{move},k}/N_p$ in P1F-B.\\
$a_i$ & Pre-transition parent index represented by output particle $i$ after a synchronous ACO transition.\\
$f_{\mathrm{parent},k}$ & Unique-parent fraction after the ACO transition, equal to the number of distinct represented pre-transition parents divided by $N_p$.\\
$m_{\max,k}$ & Maximum destination multiplicity: largest number of moving source particles that copy the same destination at sample $k$.\\
$f_{\mathrm{correct}}$ & Fraction of P1F-C evaluation runs satisfying the terminal correct-mode-lock criterion.\\
$f_{\mathrm{wrong}}$ & Fraction of P1F-C evaluation runs satisfying the terminal wrong-mode-lock criterion.\\
$f_{\mathrm{collapse}}$ & Fraction of P1F-C evaluation runs in which the unique-parent fraction drops below the catastrophic-collapse threshold at least once.\\
\bottomrule
\end{longtable}
